\begin{sol}[Dos masas en varillas paralelas unidas por un resorte]

\subsubsection*{Pregunta 1: Posiciones y velocidades}

Con el extremo izquierdo de la varilla inferior como origen $(0,0)$, las posiciones de las
masas $Q$ y $P$ son
$$\vec{r}_Q = x_1\,\hat{\imath},\qquad \vec{r}_P = x_2\,\hat{\imath} + \ell_0\,\hat{\jmath},$$
y derivando respecto al tiempo,
$$\vec{v}_Q = \dot{x}_1\,\hat{\imath},\qquad \vec{v}_P = \dot{x}_2\,\hat{\imath}.$$

\subsubsection*{Pregunta 2: Lagrangiano}

La energía cinética total es la suma de las traslaciones de ambas masas:
$$T = \tfrac12 m\dot{x}_1^2 + \tfrac12 m\dot{x}_2^2.$$

La energía potencial combina el término gravitatorio (constante, pues $P$ permanece a
altura $\ell_0$) y el elástico. Como el resorte tiene longitud natural nula, su energía es
$\tfrac12 K d^2$ con $d^2=(x_2-x_1)^2+\ell_0^2$ (teorema de Pitágoras):
$$V = m g\ell_0 + \tfrac12 K\big[(x_2-x_1)^2 + \ell_0^2\big].$$

El Lagrangiano $\mathcal{L}=T-V$ resulta
$$\boxed{\;\mathcal{L} = \tfrac12 m(\dot{x}_1^2 + \dot{x}_2^2)
- m g\ell_0 - \tfrac12 K\big[(x_2-x_1)^2 + \ell_0^2\big]\;}$$
Los términos $m g\ell_0$ y $\tfrac12 K\ell_0^2$ son constantes y no contribuyen a las
ecuaciones de movimiento.

\subsubsection*{Pregunta 3: Ecuaciones de movimiento}

Aplicando Euler-Lagrange a cada coordenada:
$$m\ddot{x}_1 = K(x_2-x_1),\qquad m\ddot{x}_2 = -K(x_2-x_1),$$
es decir,
$$\boxed{\begin{cases} m\ddot{x}_1 + K(x_1-x_2) = 0 \\[2pt] m\ddot{x}_2 + K(x_2-x_1) = 0\end{cases}}$$

Sumando ambas se obtiene $m(\ddot{x}_1+\ddot{x}_2)=0$: el centro de masa se mueve con
velocidad constante. Restándolas, la coordenada relativa $u=x_2-x_1$ satisface
$\ddot{u} + \tfrac{2K}{m}u = 0$, un movimiento armónico simple de frecuencia
$\omega=\sqrt{2K/m}$.

\end{sol}
